\begin{proposition}[The integral topology of the pinching]\hypertarget{tag-01R4}{}\label{fam:src:E.1.4}\label{fam:st:E.1.4}
\StatementRecord{fam:src:E.1.4}{proved}{01R4}{\SeriesRef{fam:src:E.1.1}; \cite[Lemma 2.1]{Engel2026}}

The compact pinching $Q$ of Proposition~\ref{fam:src:E.1.1} has
\[
 \pi_1(Q)=\Z,\qquad
 (H_i(Q;\Z))_{i=0}^6
   =(\Z,\Z,\Z^8,\Z^7,\Z^{10},0,\Z),\qquad\chi(Q)=12.
\]

\end{proposition}



\begin{proof}
\begin{enumerate}[label=\arabic*.,leftmargin=*]
\item The two section embeddings are cofibrations: a tubular neighbourhood
retracts onto each section, and the two neighbourhoods are disjoint.
The ordinary pushout is therefore homotopy equivalent to the space
formed from $R$ by attaching $S\times[0,1]$ along
$s_0$ and $s_\infty T$. This comparison follows by collapsing the
intervals in the mapping cylinder; the tubular retractions construct
its homotopy inverse. Both $S$ and $R$ are simply connected. Van
Kampen applied to this mapping-cylinder model gives one free loop
and no additional generator or relation, so $\pi_1(Q)=\Z$.
The same model has the exact sequence
\begin{equation}
 H_i(S)\xrightarrow{d_i=s_{0*}-s_{\infty*}T_*}H_i(R)
 \longrightarrow H_i(Q)\longrightarrow H_{i-1}(S).
\label{pub:I:eq:E4}
\end{equation}

\item Put $F=[S_t]$, $m=P-O$, and
$v_0=[6P]-[O]-3F\in H_2(S;\Z)$. The intersection computations in
\cite[Lemma~2.1, pp.~3--4]{Engel2026} give
\label{layout:I:2}
\begin{equation}
\begin{gathered}
F^2=Fv_0=Fm=0,\quad v_0^2=-6,\quad mv_0=-1,\quad m^2=-2,
 \\
\quad OF=1,\quad Ov_0=0,\quad Om=1.
\end{gathered}
\label{pub:I:eq:E5}
\end{equation}
The translation preserves all components of the two reducible
Kodaira fibres. Their root classes span a negative definite space of
rank seven. Its orthogonal complement is generated over $\Q$ by
$F,O,v_0$. Since $T_*O=O+v_0+3F$, preservation of the intersection
form gives, on this complement and hence on the entire lattice,
\begin{equation}
 T_*a=a+(aF)v_0-\bigl(av_0-3aF\bigr)F.
\label{pub:I:eq:E6}
\end{equation}
Its value on $v_0$ is $v_0+6F$: it has no $O$-component
and has zero intersections with the rank-seven root space, so it is
$\lambda v_0+aF$. The already established identity $T^*m=m-F$ gives
$mT_*v_0=(m-F)v_0=-1$, forcing $\lambda=1$. Orthogonality to
$T_*O$ then forces $a=6$. These
values on $F,O,v_0$ and the root space verify \eqref{pub:I:eq:E6} on a rational basis,
and both sides preserve the integral lattice. In particular
$T^*m=m-F$, in agreement with \ref{fam:src:E.1.1}.

\item The rational elliptic surface is the blow-up of $\mathbb P^2$ at nine
points, so its homology in even degrees is $\Z,\Z^{10},\Z$ and
vanishes in odd degrees; its intersection lattice is unimodular.
The projective-bundle theorem follows here by local trivialization
and the fibre hyperplane class: $H_2(R)=H_2(S)\oplus\Z f$ and
$H_4(R)=\Z\oplus H_2(S)$, with $f$ the projective-line fibre.
The hyperplane class that vanishes on $s_0$ restricts to $-m$ on
$s_\infty$. Therefore
\begin{equation}
 d_2(a)=(1-T_*)a+(mT_*a)f
       =-(aF)v_0+(av_0-3aF)F+(am-aF)f.
\label{pub:I:eq:E7}
\end{equation}
The span of $F,v_0,m$ is primitive. To check the only nontrivial
part of this assertion, split off the unimodular plane spanned by
$F,O$; $m-F$ and $v_0$ lie in its negative definite orthogonal
complement. Their Gram determinant is
$\det\left(\begin{smallmatrix}-2&-1\\-1&-6\end{smallmatrix}\right)=11$.
The square of an index of a larger integral lattice would divide
eleven, so the index is one. The span of $F,v_0$ is also primitive:
pairing with $O,m$ gives determinant $-1$. Unimodularity consequently
makes $a\mapsto(aF,av_0,am)$ surjective onto $\Z^3$. The coefficient
matrix in \eqref{pub:I:eq:E7} has determinant $-1$. Thus $d_2$ has primitive image
$\Z v_0\oplus\Z F\oplus\Z f$, rank three, kernel $\Z^7$,
and cokernel $\Z^8$.

\item The map $d_4$ sends $[S]$ to $s_0[S]-s_\infty[S]$, whose Poincar\'e
dual is $\pi_R^*m$. This is primitive and nonzero since $mO=1$.
Thus $d_4$ is injective with cokernel $\Z^{10}$. The map $d_0$ is
zero. Substitution in \eqref{pub:I:eq:E4} proves every asserted homology group.
The Euler characteristic is independently obtained by stratification:
$\chi(R)-\chi(S)=24-12=12$. 
\end{enumerate}
\end{proof}